\section{Local counts and cluster inversion}\label{inf:sec-localization}

The window in this section is an arbitrary fixed number
$W_{\mathrm{loc}}>0$.
The local count depends on this number, but neither on the centre
order nor on the detuning or split.

\begin{theorem}[Local logarithmic count]\label{inf:local-count}
There is an integer $M_{W_{\mathrm{loc}}}\ge1$ with the following
quantifier order: for every fixed $D>0$ there is
$P_{W_{\mathrm{loc}},D}<\infty$ such that every angular interval
of length $D\log p$ contains at most $M_{W_{\mathrm{loc}}}$
states of $\mathcal W_p(W_{\mathrm{loc}})$ whenever
$p\ge P_{W_{\mathrm{loc}},D}$ and $R\in[2p-1,2p+4]$.
In particular the count bound is independent of $D$.
\end{theorem}
\begin{proof}
In the oscillatory region $R-\nu\ge p^{2/3}/2$,
Lemma~\ref{inf:phase-strip} assigns to each state a lattice point
with $|m-\varphi_p(j)|\le C_{W_{\mathrm{loc}}}/p$.
Dividing \eqref{inf:phase-strip-bound} by $\pi$ gives the
stronger bound $C_{W_{\mathrm{loc}}}/(\pi p)$; we use the same
constant in the weaker bound above. The threshold includes those
of Lemmas~\ref{inf:window-isolation} and \ref{inf:phase-strip}
for this window and the condition $2C_{W_{\mathrm{loc}}}/p<1$.
Its curvature satisfies
\[
 c/p\le\varphi_p''(j)\le Cp^{-5/6}.
\]
The lower bound uses $\kappa\le Cp$; the upper uses
$\kappa\ge cp^{5/6}$. On an interval of length $D\log p$,
\[
 \tfrac12Cp^{-5/6}(D\log p)^3
       +2C_{W_{\mathrm{loc}}}D\log p/p<1
\]
for large $p$. Lemma~\ref{inf:lattice-strip} therefore bounds the
number of states by
$2+4\sqrt{C_{W_{\mathrm{loc}}}/c}$.
This number is independent of $D$; only the threshold of the
last inequality depends on $D$.

In the turning region $\nu\ge R-2p^{2/3}$, suppose two states
have angular distance at most $D\log p$.
Their original roots are within $O_{W_{\mathrm{loc}}}(p^{-1})$
of $R$. Transport the root of the larger order to the smaller
order $\nu_0$, keeping its radial index fixed.
Lemma~\ref{inf:bessel-order} moves it by at most
$\pi D\log p$. Every intervening root stays of order $p$.
If the two radial indices were equal, the lower derivative bound
in that lemma would make the original roots differ by at least
$2(1-o(1))$, because the angular indices differ by at least one.
That contradicts their original window width $O(p^{-1})$.
Thus there are two distinct zeros of $J_{\nu_0}$ in an interval
of length $O_{W_{\mathrm{loc}},D}(\log p)$ about $R$.

Throughout that interval the coefficient in
\[
 v''+\left(1-\frac{\nu_0^2-1/4}{x^2}\right)v=0,
 \qquad v(x)=\sqrt xJ_{\nu_0}(x),
\]
is at most $(2+o(1))p^{-1/3}$, uniformly for the fixed window
and interval constant, because
$\nu_0\ge R-2p^{2/3}$ and $x=R+O(\log p)$.
Between the two zeros, integration and the Dirichlet Poincar\'e
inequality give
$\pi^2/\ell^2\le Cp^{-1/3}$ for their separation $\ell$.
Hence $\ell\ge(\pi/\sqrt2+o(1))p^{1/6}$, in particular
$\ell\ge2.2p^{1/6}$ for large $p$, contradicting its logarithmic upper
bound for large $p$. The turning interval therefore contains
at most one state.

The two regions overlap by an order interval of length
$3p^{2/3}/2$. An angular interval of length $D\log p$ has order
length $2D\log p$, which is smaller than this overlap for large
$p$. If it is not entirely in the oscillatory region, all its
possible window points are in the turning region.
There is at most one radial state per angular index by
Lemma~\ref{inf:window-isolation}. Rounding up the oscillatory
count and increasing it to at least one proves the theorem.
\end{proof}

\begin{corollary}[Bounded cluster cardinality]\label{inf:clusters}
Fix $D>0$ and connect consecutive angular indices of
$\mathcal W_p(W_{\mathrm{loc}})$ when their separation is at most
$D\log p$. For all sufficiently large $p$, every resulting
cluster has at most $M_{W_{\mathrm{loc}}}$ states, and distinct
clusters are separated by more than $D\log p$.
A cluster may have diameter as large as
$(M_{W_{\mathrm{loc}}}-1)D\log p$.
The same cardinality bound holds for every subset of the window.
\end{corollary}
\begin{proof}
If a cluster had $M_{W_{\mathrm{loc}}}+1$ states, its first that
many states would span at most
$M_{W_{\mathrm{loc}}}D\log p$.
Apply Theorem~\ref{inf:local-count} with the fixed constant
$M_{W_{\mathrm{loc}}}D$ in place of $D$.
Its bound is still $M_{W_{\mathrm{loc}}}$, a contradiction.
The separation and diameter statements follow from the definition.
Removing states can split clusters but cannot join two clusters
of the larger set, proving the assertion for subsets.
\end{proof}

\begin{lemma}[Inversion with arbitrary diagonal Hilbert weights]
\label{inf:cluster-inverse}
Fix positive constants $c,C,q_{\mathrm{gap}},\sigma$ and an integer $M\ge1$.
Let $D$ satisfy $\sigma D>q_{\mathrm{gap}}M+q_{\mathrm{gap}}+2$.
Let $\tau>0$ and let $F$ be a finite matrix whose indices
carry angular degrees, partitioned into clusters of at most $M$
indices separated by more than $D\log p$.
Assume that:
\begin{enumerate}[label=\textup{(\roman*)}]
\item $F$ is self-adjoint for a positive diagonal Hilbert inner
product and $\nrm{Fv}_H\ge\gamma_F\nrm v_H$, where
$\gamma_F=c\tau p^{-q_{\mathrm{gap}}}$;
\item in a weighted coefficient $\ell^1$ norm,
$\nrm F_1\le C\tau$;
\item deleting all inter-cluster entries has error at most
$C\tau p^{-\sigma D}$ in both that coefficient norm and the
specified Hilbert operator norm.
\end{enumerate}
Then, for sufficiently large $p$ depending only on
$c,C,q_{\mathrm{gap}},\sigma,M,D$,
\begin{equation}\label{inf:cofactor-inverse-bound}
 \nrm{F^{-1}}_1\le C'\tau^{-1}p^{q_{\mathrm{gap}}M}.
\end{equation}
The assertion is uniform as $\tau\downarrow0$.
\end{lemma}
\begin{proof}
Let $F_{\mathrm b}$ keep the diagonal cluster blocks.
The symmetric deletion of paired entries preserves
self-adjointness in the diagonal Hilbert structure.
The Hilbert error is at most $\gamma_F/2$ for large $p$, since
$\sigma D>q_{\mathrm{gap}}$. Thus
\[
 \nrm{F_{\mathrm b}v}_H\ge
 (\gamma_F-C\tau p^{-\sigma D})\nrm v_H\ge(\gamma_F/2)\nrm v_H.
\]
Every eigenvalue of every cluster block is real and has modulus
at least $\gamma_F/2$.

Conjugate a cluster block of size $m\le M$ by its coefficient
weights, so its $\ell^1$ norm is the unweighted column norm.
All its entries have modulus at most $C\tau$.
The determinant is invariant under this conjugation and has
modulus at least $(\gamma_F/2)^m$.
Expansion of each cofactor as a sum of $(m-1)!$ products gives
its absolute value at most $(m-1)!(C\tau)^{m-1}$.
For $m=1$ the empty product is one.
Cramer's formula and summing each inverse column give
\[
 \nrm{F_{\mathrm b,m}^{-1}}_1
 \le m!(C\tau)^{m-1}(2/\gamma_F)^m
 \le C_{M,c,C}\tau^{-1}p^{q_{\mathrm{gap}}M}.
\]
The block diagonal inverse has the maximum of these norms.
Its product with the discarded coefficient error is at most
$C'p^{q_{\mathrm{gap}}M-\sigma D}$, below $1/2$ for large $p$.
A Neumann series restores that error and proves
\eqref{inf:cofactor-inverse-bound}.
In both absorptions the factor $\tau$ cancels. The inverse has
exactly one remaining factor $\tau^{-1}$.
For an empty index set the bound is vacuous.
\end{proof}

The next section uses a window scaled by $\tau$,
and applies the local count to a fixed larger window containing it.
